\section{Strumenti utilizzati}

Di seguito vengono illustrati gli strumenti e le tecnologie adottate per la realizzazione
del progetto. La scelta dello stack tecnologico è stata orientata a garantire la massima
integrazione tra la gestione dei dati geografici e un’architettura di tipo modulare.

\begin{itemize}

\item \textbf{Spring Boot:} Framework Java basato sull’ecosistema Spring, impiegato per implementare
la logica di backend secondo il pattern architetturale MVC. Spring Boot
facilita la gestione del Controller e l’integrazione del Model con i database, automatizzando
la configurazione dell’infrastruttura. Questo approccio garantisce una
chiara separazione delle responsabilità, rendendo l’applicazione scalabile, modulare
e facilmente manutenibile.

\item \textbf{GraphHopper:} motore di routing utilizzato per ottenere dati reali di distanza e
tempo di percorrenza tra le coordinate GPS dei punti di interesse. Fornisce la
componente pratica necessaria per calcolare percorsi realistici sulle strade esistenti,
garantendo itinerari ottimizzati basati su tempi reali di percorrenza e condizioni
stradali.

\item \textbf{JUnit 4:} Framework per il testing unitario in Java, impiegato per validare il corretto
funzionamento delle classi e dei metodi sviluppati. L’uso di test automatizzati
consente di rilevare tempestivamente eventuali errori e di garantire la robustezza del
codice.

\item \textbf{MongoDB e MongoDB Compass:} MongoDB è il database NoSQL utilizzato per
memorizzare i POI (Points of Interest) con strutture flessibili e attributi variabili
provenienti da OpenStreetMap, come nome, descrizione, coordinate, categorie e
immagini. Spring Data MongoDB facilita le operazioni CRUD e le query complesse.
MongoDB Compass, invece, è lo strumento grafico che permette di esplorare
visivamente le collezioni, verificare la correttezza dei documenti JSON e correggere
manualmente eventuali dati errati, migliorando il controllo sulla qualità dei dati.

\item \textbf{OpenStreetMap:} piattaforma open-source che fornisce dati geospaziali dettagliati.
I dati di OpenStreetMap sono utilizzati da Nominatim per il geocoding, da Overpass
API per interrogazioni sui monumenti e da React-Leaflet per la visualizzazione
delle mappe interattive nel frontend. La natura collaborativa di OpenStreetMap
garantisce dati aggiornati e accurati per il calcolo dei percorsi.

\item \textbf{Nominatim:} servizio di geocoding utilizzato per trasformare indirizzi testuali inseriti
dall’utente in coordinate geografiche precise ($lat/lon$), consentendo di posizionare
correttamente il punto di partenza dei percorsi sulla mappa.

\item \textbf{Overpass API:} interroga il database OpenStreetMap per individuare monumenti e
POI nelle vicinanze del punto di interesse indicato dall’utente. Per evitare di sovraccaricare
il servizio e garantire tempi di risposta rapidi, i dati vengono memorizzati
localmente su MongoDB dopo la prima richiesta, rendendo più efficienti le ricerche
successive.

\item \textbf{React.js \& Node.js:} React gestisce lo stato dell’interfaccia utente, come la lista dei
monumenti selezionati e le informazioni sul percorso ottimizzato, mentre Node.js
offre l’ambiente di esecuzione e la gestione dei pacchetti tramite npm. Insieme,
permettono di costruire un frontend dinamico, reattivo e facilmente aggiornabile.

\item \textbf{React-Leaflet:} libreria che connette React alle mappe interattive di OpenStreetMap.
Consente di posizionare marker sui monumenti, visualizzare linee di percorso ottimizzate
calcolate dal backend e rendere l’esperienza dell’utente più intuitiva grazie
a mappe dinamiche e interattive.

\item \textbf{Axios:} libreria per gestire chiamate HTTP asincrone. Permette al frontend di inviare
richieste al backend e ricevere le risposte in formato JSON senza ricaricare la pagina,
garantendo un’esperienza utente fluida e immediata.

\item \textbf{Eclipse:} IDE principale per la scrittura del codice Java, il debugging del backend e
la gestione del progetto tramite Maven o Gradle. Offre strumenti avanzati di refactoring,
analisi del codice e integrazione con sistemi di versionamento, migliorando
l’efficienza dello sviluppo.

\item \textbf{CodeMR:} strumento di analisi statica del codice sorgente, che estrae metriche di
coesione e accoppiamento delle classi. Aiuta a identificare classi troppo complesse
e suggerisce possibili rifattorizzazioni, migliorando la manutenibilità del software.

\item \textbf{Postman:} Strumento essenziale per il testing delle API REST sviluppate con Spring
Boot. Consente di effettuare richieste HTTP, validare le risposte del server e automatizzare
test, facilitando così il debug e l’integrazione tra i diversi componenti
del sistema.

\item \textbf{StarUML \& Diagrams.net:} strumenti utilizzati per la modellazione dei sistemi.
StarUML serve per diagrammi formali (casi d’uso, classi, deployment), mentre
Diagrams.net è utile per schemi logici, flussi di dati e rappresentazioni più informali,
favorendo la chiarezza della progettazione.

\item \textbf{GitHub \& GitHub Desktop:} strumenti di controllo di versione che permettono
di salvare lo storico dei file, gestire branch multipli e sincronizzare il lavoro tra
diversi membri del team, garantendo una collaborazione efficiente e tracciabilità
delle modifiche.

\item \textbf{WhatsApp:} applicazione di messaggistica utilizzata come strumento di comunicazione
interna tra i membri del team. Consente di coordinare le attività di sviluppo,
discutere problemi tecnici e organizzare riunioni in tempo reale, garantendo un
flusso comunicativo rapido ed efficiente.

\end{itemize}